\documentclass[10pt,a4paper]{amsart}
\usepackage{amsmath,amssymb,amsthm,graphicx}
% Retained preamble for the Stacks Project chapter snapshot in this directory.
% The chapter's own preamble.tex is a copy of the shared Stacks Project preamble; the block
% below is sliced MECHANICALLY, line for line, from it. Nothing was added, removed or
% redefined; the wrapper class, the xr-hyper/hyperref wiring and the externaldocument list
% are omitted because the assembler supplies the document class.
%
% lines taken from preamble.tex: 3, 5, 9, 11, 12, 156, 157, 158, ... (51 lines total)
%
% package coverage restored from the chapter's own preamble (2026-10-13): xy, xy-2cell

\usepackage[all]{xy}
\xyoption{2cell}
\UseAllTwocells

\newenvironment{reference}{\comment}{\endcomment}
\newenvironment{slogan}{\comment}{\endcomment}
\newenvironment{history}{\comment}{\endcomment}
\theoremstyle{plain}
\newtheorem{theorem}[subsection]{Theorem}
\newtheorem{proposition}[subsection]{Proposition}
\newtheorem{lemma}[subsection]{Lemma}
\theoremstyle{definition}
\newtheorem{definition}[subsection]{Definition}
\newtheorem{example}[subsection]{Example}
\newtheorem{exercise}[subsection]{Exercise}
\newtheorem{situation}[subsection]{Situation}
\theoremstyle{remark}
\newtheorem{remark}[subsection]{Remark}
\newtheorem{remarks}[subsection]{Remarks}
\numberwithin{equation}{subsection}
\def\lim{\mathop{\mathrm{lim}}\nolimits}
\def\colim{\mathop{\mathrm{colim}}\nolimits}
\def\Spec{\mathop{\mathrm{Spec}}}
\def\Hom{\mathop{\mathrm{Hom}}\nolimits}
\def\Ext{\mathop{\mathrm{Ext}}\nolimits}
\def\SheafHom{\mathop{\mathcal{H}\!\mathit{om}}\nolimits}
\def\SheafExt{\mathop{\mathcal{E}\!\mathit{xt}}\nolimits}
\def\Sch{\mathit{Sch}}
\def\Mor{\mathop{\mathrm{Mor}}\nolimits}
\def\Ob{\mathop{\mathrm{Ob}}\nolimits}
\def\Sh{\mathop{\mathit{Sh}}\nolimits}
\def\NL{\mathop{N\!L}\nolimits}
\def\CH{\mathop{\mathrm{CH}}\nolimits}
\def\proetale{{pro\text{-}\acute{e}tale}}
\def\etale{{\acute{e}tale}}
\def\QCoh{\mathit{QCoh}}
\def\Ker{\mathop{\mathrm{Ker}}}
\def\Im{\mathop{\mathrm{Im}}}
\def\Coker{\mathop{\mathrm{Coker}}}
\def\Coim{\mathop{\mathrm{Coim}}}
\DeclareMathSymbol{\boxtimes}{\mathbin}{AMSa}{"02}
\def\QCohstack{\mathcal{QC}\!\mathit{oh}}
\def\Cohstack{\mathcal{C}\!\mathit{oh}}
\def\Spacesstack{\mathcal{S}\!\mathit{paces}}
\def\Quotfunctor{\mathrm{Quot}}
\def\Hilbfunctor{\mathrm{Hilb}}
\def\Curvesstack{\mathcal{C}\!\mathit{urves}}
\def\Polarizedstack{\mathcal{P}\!\mathit{olarized}}
\def\Complexesstack{\mathcal{C}\!\mathit{omplexes}}
\def\Pic{\mathop{\mathrm{Pic}}\nolimits}
\def\Picardstack{\mathcal{P}\!\mathit{ic}}
\def\Picardfunctor{\mathrm{Pic}}
\def\Deformationcategory{\mathcal{D}\!\mathit{ef}}
\usepackage{cleveref}
\crefname{originallemma}{original lemma}{original lemmas}
\Crefname{originallemma}{Original lemma}{Original lemmas}
\makeatletter
\newlabel{algebra-lemma-simple-pieces}{{00J3}{source}{}{sourcepointer-1}{}}
\newlabel{algebra-lemma-simple-pieces@cref}{{[originallemma][0][]00J3}{[1][0][]source}{}{}{}}
\makeatother
\crefname{originallemma}{original lemma}{original lemmas}
\Crefname{originallemma}{Original lemma}{Original lemmas}
\makeatletter
\newlabel{fields-lemma-separable-trace-pairing}{{0BIL}{source}{}{sourcepointer-2}{}}
\newlabel{fields-lemma-separable-trace-pairing@cref}{{[originallemma][0][]0BIL}{[1][0][]source}{}{}{}}
\makeatother
\crefname{originallemma}{original lemma}{original lemmas}
\Crefname{originallemma}{Original lemma}{Original lemmas}
\makeatletter
\newlabel{lemma-finite-flat-trace}{{0BT8}{source}{}{sourcepointer-3}{}}
\newlabel{lemma-finite-flat-trace@cref}{{[originallemma][0][]0BT8}{[1][0][]source}{}{}{}}
\makeatother
\title{Stacks, Discriminants: lemma-tau-nonzero}
\author{The Stacks Project Authors}
\date{stacks-project:discriminant.tex:local-snapshot-20261012; GNU FDL 1.2 or later}
\begin{document}
\maketitle
\noindent\textit{Editorial note.} Statement span is lines 1126-1136 of sources/stacks-discriminant/pinned-original-source.tex; proof span is lines 1138-1159, the complete author proof (20 body lines).\par\medskip
\noindent\textit{Editorial note.} Navigation stubs: 3 cited-but-not-included labels are supplied as non-load-bearing pointers with their pinned tags where the snapshot tag index has them.\par\medskip
\noindent\textit{Editorial note.} extra\_packages is empty on purpose: the retained preamble already loads xy with options and a second declaration clashes in the assembler wrapper.\par\medskip
\noindent\textit{Editorial note.} Text is retained verbatim, including the author's own typos and placeholders.\par\medskip
\noindent\textit{Editorial note.} Queue index 30 of queue\_discriminant.json (pb=882, stmt 1126-1136, proof 1138-1159).\par\medskip
\noindent\textit{Allowed standard prerequisites (editorial).}\begin{itemize}
\item Explicit prerequisite boundary. The cited inputs are prerequisites; the argument in the retained span is not replaced by any cited result.
\end{itemize}\medskip
\ifdefined\hypertarget\hypertarget{sourcepointer-1}{}\fi
\noindent\textit{Source reference.} Original lemma 00J3: sources/stacks-discriminant/pinned-original-source.tex, source lines 1159--1159. Cited by the retained statement or proof but not part of the collected spans; supplied only as a navigation pointer. The argument does not depend on its content.\par\medskip
\ifdefined\hypertarget\hypertarget{sourcepointer-2}{}\fi
\noindent\textit{Source reference.} Original lemma 0BIL: sources/stacks-discriminant/pinned-original-source.tex, source lines 1159--1159. Cited by the retained statement or proof but not part of the collected spans; supplied only as a navigation pointer. The argument does not depend on its content.\par\medskip
\ifdefined\hypertarget\hypertarget{sourcepointer-3}{}\fi
\noindent\textit{Source reference.} Original lemma 0BT8: sources/stacks-discriminant/pinned-original-source.tex, source lines 874--877. Cited by the retained statement or proof but not part of the collected spans; supplied only as a navigation pointer. The argument does not depend on its content.\par\medskip
\begin{lemma}
\label{lemma-tau-nonzero}
Let $k$ be a field and let $A$ be a finite $k$-algebra. Assume $A$
is local with residue field $k'$. The following are equivalent
\begin{enumerate}
\item $\text{Trace}_{A/k}$ is nonzero,
\item $\tau_{A/k} \in \omega_{A/k}$ is nonzero, and
\item $k'/k$ is separable and $\text{length}_A(A)$ is prime
to the characteristic of $k$.
\end{enumerate}
\end{lemma}

\begin{proof}
Conditions (1) and (2) are equivalent by Lemma \ref{lemma-finite-flat-trace}.
Let $\mathfrak m \subset A$. Since $\dim_k(A) < \infty$ it is clear that
$A$ has finite length over $A$. Choose a filtration
$$
A = I_0 \supset \mathfrak m = I_1 \supset I_2 \supset \ldots I_n = 0
$$
by ideals such that $I_i/I_{i + 1} \cong k'$ as $A$-modules. See
Algebra, Lemma \ref{algebra-lemma-simple-pieces} which also shows that
$n = \text{length}_A(A)$. If $a \in \mathfrak m$ then $aI_i \subset I_{i + 1}$
and it is immediate that $\text{Trace}_{A/k}(a) = 0$.
If $a \not \in \mathfrak m$ with image $\lambda \in k'$, then
we conclude
$$
\text{Trace}_{A/k}(a) =
\sum\nolimits_{i = 0, \ldots, n - 1}
\text{Trace}_k(a : I_i/I_{i - 1} \to I_i/I_{i - 1}) =
n \text{Trace}_{k'/k}(\lambda)
$$
The proof of the lemma is finished by applying
Fields, Lemma \ref{fields-lemma-separable-trace-pairing}.
\end{proof}

\end{document}
